\chapter{Exact Cancellation in a Normalized Two-Channel State}
\label{chap:cancellation}

\section{The complementary amplitude}

Let $\mathcal{H}=\C^2$ with orthonormal basis $\{\ket{0},\ket{1}\}$. For $0\le\sigma\le1$ and phase $\phi\in\R$, define the normalized state
\begin{equation}
 \ket{\psi(\sigma,\phi)}
 =\sqrt{\sigma}\ket{0}+e^{i\phi}\sqrt{1-\sigma}\ket{1}.
 \label{eq:binary-state}
\end{equation}
Let
\[
 \ket{+}=\frac{\ket{0}+\ket{1}}{\sqrt2},
 \qquad
 \ket{-}=\frac{\ket{0}-\ket{1}}{\sqrt2}.
\]
The unnormalized equal-gain readout amplitude is
\begin{equation}
 \Acomp(\sigma,\phi)
 =\sqrt{\sigma}+e^{i\phi}\sqrt{1-\sigma}
 =\sqrt2\,\braket{+}{\psi(\sigma,\phi)}.
 \label{eq:Acomp-def}
\end{equation}
Thus exact cancellation means that the state is orthogonal to the symmetric detector $\ket{+}$.

\section{Uniqueness theorem}

Squaring the modulus gives
\begin{equation}
 |\Acomp(\sigma,\phi)|^2
 =1+2\sqrt{\sigma(1-\sigma)}\cos\phi.
 \label{eq:Acomp-square}
\end{equation}

\begin{lemma}[SOH-L002: exact complementary cancellation]
\label{lem:exact-cancellation}
For $0\le\sigma\le1$,
\[
 \Acomp(\sigma,\phi)=0
 \quad\Longleftrightarrow\quad
 \sigma=\frac12,
 \quad \phi\equiv\pi\pmod{2\pi}.
\]
\end{lemma}
\begin{proof}
If $\Acomp=0$, then
\[
 \sqrt{\sigma}=-e^{i\phi}\sqrt{1-\sigma}.
\]
Taking moduli gives $\sigma=1-\sigma$, hence $\sigma=1/2$. Substituting back yields $e^{i\phi}=-1$. The converse follows immediately.
\end{proof}

The proof contains two logically distinct requirements:
\begin{enumerate}
  \item equal magnitudes force balance, $\sigma=1/2$;
  \item opposite phase forces $\phi=\pi$ modulo $2\pi$.
\end{enumerate}
Neither condition alone is sufficient.

\begin{figure}[htbp]
 \centering
 \includegraphics[width=0.82\textwidth]{cancellation_landscape.png}
 \caption{The equal-gain cancellation landscape has a unique zero at $(\sigma,\phi)=(1/2,\pi)$ modulo phase periodicity.}
 \label{fig:cancellation-landscape}
\end{figure}

\section{Geometric proof and the triangle inequality}

The same result follows from the equality condition in the triangle inequality. Two complex numbers of lengths $\sqrt\sigma$ and $\sqrt{1-\sigma}$ can sum to zero only if their lengths are equal and their directions are opposite. This geometric formulation is valuable because it generalizes to Hilbert-space components: exact cancellation against a rank-one detector is an equality case of a norm inequality.

At the unique zero,
\[
 \ket{\psi(1/2,\pi)}=\ket{-},
\]
up to a global phase. Therefore the cancellation condition selects not merely the equator of the Bloch sphere but one ray on that equator.

\section{Local stability}

Set
\[
 \sigma=\frac12+x,
 \qquad
 \phi=\pi+\delta.
\]
Then
\begin{align}
 |\Acomp|^2
 &=1-\sqrt{1-4x^2}\cos\delta\notag\\
 &=2x^2+\frac12\delta^2+O\!\left(x^4+x^2\delta^2+\delta^4\right).
 \label{eq:cancellation-local}
\end{align}
The cancellation point is therefore a non-degenerate quadratic minimum in the two transverse coordinates. Horizontal imbalance costs $2x^2$ to leading order, the same coefficient that appears in the entropy deficit \cref{eq:entropy-deficit-series}. Phase mismatch costs $\delta^2/2$.

This agreement of the horizontal quadratic terms is not accidental. Both arise from a symmetric binary normalization near equal weights. It provides a precise local sense in which information imbalance and failure of destructive interference measure the same displacement. It does not identify either with $|\xi(s)|$.

\begin{figure}[htbp]
 \centering
 \includegraphics[width=0.80\textwidth]{cancellation_quadratic_error.png}
 \caption{Error of the quadratic approximation in \cref{eq:cancellation-local}. The exact minimum is locally coercive.}
 \label{fig:cancellation-quadratic}
\end{figure}

\section{Weighted detectors and the hidden equal-gain assumption}

The half point depends on equal channel gain. Let $a,b>0$ and define
\begin{equation}
 \Acomp_{a,b}(\sigma,\phi)
 =a\sqrt\sigma+b e^{i\phi}\sqrt{1-\sigma}.
 \label{eq:weighted-amplitude}
\end{equation}
Then exact cancellation requires
\begin{equation}
 \sigma=\frac{b^2}{a^2+b^2},
 \qquad \phi\equiv\pi\pmod{2\pi}.
 \label{eq:weighted-balance}
\end{equation}
The balance point is $1/2$ if and only if $a=b$.

This fact is central to claim discipline. The zeta-state ansatz cannot simply choose an equal-gain detector for aesthetic reasons. Equal gain must be justified by involution symmetry, unitarity, a canonical inner product, or another independently derived principle. Otherwise a hidden metric asymmetry shifts the selected line.

\section{Why two channels matter}

For three or more channels, exact vector cancellation can occur for many unequal weight configurations. The uniqueness of $1/2$ is a special property of a complementary pair with a fixed equal-gain readout. Thus the word ``binary'' is mathematically substantive. A many-term alternating Dirichlet series does not automatically reduce to the two-channel lemma; a canonical coarse-graining must be constructed.

\section{Readout factorization target}

The most direct bridge would have the form
\begin{equation}
 \xi(s)=g(s)\,\Acomp\bigl(\Re s,\Phi(s)\bigr),
 \label{eq:naive-factorization}
\end{equation}
where $g$ never vanishes in $\Sstrip$. Then any zero of $\xi$ would be a zero of $\Acomp$, and \cref{lem:exact-cancellation} would prove RH.

However, \cref{eq:naive-factorization} is not yet a legitimate construction. The right-hand amplitude depends explicitly on $\Re s$ and is generally non-holomorphic; $\Phi$ and $g$ can be engineered to reproduce any zero set; normalization may become singular; and no canonical origin has been supplied. These obstructions are treated in \cref{chap:nogo,chap:candidates}.

\status{SOH-L002, the local expansion, and the weighted-channel result are exact. Their use as a zeta zero detector is conditional on a canonical readout factorization.}
